\documentclass[11pt,a4paper]{article}

\usepackage{../shared/cathedral-preamble}
\usepackage{setspace}
\onehalfspacing

\title{%
  An Invitation to Complete the Cathedral:\\
  Two Axioms Remain%
}
\author{
  Jason Robert Gochanour
}
\date{April 26, 2026}

\begin{document}

\maketitle

\begin{center}
\emph{For the student who hasn't started yet.}
\end{center}

\vspace{1cm}

% ================================================================
\section*{What This Is}
% ================================================================

We built something called the Cathedral: a 161-file computer program
that proves, with mathematical certainty, that the Riemann Hypothesis
is \emph{exactly equivalent} to the convergence of a specific
approximation problem. Every step is verified by a computer. Every
assumption is named. Every failed attempt is preserved.

The proof is complete---except for two assumptions, called
\emph{axioms}. We couldn't prove them. We named them, classified
them, and documented exactly what they say and why they're hard.

This paper is an invitation. Those two axioms are now your problem.

% ================================================================
\section*{Who Built This}
% ================================================================

A computer scientist in New Mexico. Not a professional mathematician.
No PhD in mathematics. No university affiliation. No research grant.

The tools: a laptop, an internet connection, Lean~4 (free, open
source), Mathlib (free, open source), and AI assistance (commercially
available).

The timeline: 30 days.

If you have a laptop and curiosity, you have everything that was
needed to build this. The barrier to entry for machine-verified
mathematics is now approximately zero.

% ================================================================
\section*{What You Need to Know}
% ================================================================

\subsection*{The Problem}

The Riemann Hypothesis (1859) says: all non-trivial zeros of the
Riemann zeta function $\zeta(s) = \sum_{n=1}^\infty n^{-s}$ lie
on the line $\operatorname{Re}(s) = \frac{1}{2}$.

Nobody has proved this. Nobody has disproved it. It has been open
for 167 years. There is a \$1,000,000 prize.

\subsection*{The Translation}

The Cathedral proves that this is equivalent to:

\[
  \forall\,\varepsilon > 0,\;\exists\,N_0,\;\forall\,N \geq N_0,\;
  \exists\,v \in \mathbb{R}^{N-1}:\quad
  \int_0^1 \left(1 - \sum_{k=2}^N v_k\left\{\frac{1}{kx}\right\}
  \right)^2 dx < \varepsilon.
\]

In words: you can approximate the constant function $1$ arbitrarily
well by mixing together sawtooth waves $\{1/(kx)\}$.

\subsection*{The Two Crown Axioms}

\begin{enumerate}
  \item \textbf{\texttt{critical\_line\_mellin\_variance}}: Under RH,
    the Mellin variance of the BD residual on the critical line
    decays as $O(1/\log N)$.

    \emph{What this means}: the Mellin transform of the approximation
    error, evaluated on the critical line, has bounded $L^2$ norm.
    This follows from the Hardy--Littlewood mean value theorem
    $\int_0^T |1/\zeta(1/2+it)|^2\,dt = O(T)$ (1918).
    Mathlib lacks this infrastructure.

  \item \textbf{\texttt{rh\_zeta\_lower\_bound\_from\_zero\_counting}}:
    A polynomial lower bound on $|\zeta(s)|$ away from the
    critical line.

    \emph{What this means}: the zeta function doesn't get too
    small too fast. This follows from Hadamard's product formula
    and the zero-counting function. Borel--Carath\'eodory is in
    Mathlib; the gap is connecting it to the Hadamard product
    and Riemann--von Mangoldt zero counting.
\end{enumerate}

Each axiom is independently meaningful. Proving either one
strengthens the formalization.

% ================================================================
\section*{What You Can Do Right Now}
% ================================================================

\subsection*{If You Know Some Lean~4}

\begin{enumerate}
  \item \textbf{Prove Axiom~1} (\texttt{critical\_line\_mellin\_variance}).
    This requires formalizing the Hardy--Littlewood mean value theorem
    in Lean. The approximate functional equation and fourth-moment
    method are the key tools.

  \item \textbf{Prove Axiom~2} (\texttt{rh\_zeta\_lower\_bound\_from\_zero\_counting}).
    Borel--Carath\'eodory is in Mathlib. The gap is the Hadamard
    product formula and Riemann--von Mangoldt zero counting.

  \item \textbf{Prove one of the 53 non-critical axioms.} The
    Cathedral contains 55 axioms total. Only 2 are on the crown
    path. The other 53 support alternative proof engines (spectral,
    sieve, Perron). Proving any one of them strengthens the
    overall formalization.
\end{enumerate}

\subsection*{If You Know Some Analysis}

\begin{enumerate}
  \item \textbf{Decompose Axiom~2.} The current axiom is a
    ``black box'' that packages the entire forward direction.
    A useful contribution would be to decompose it into 3--5
    intermediate lemmas, each of which is independently
    meaningful and potentially formalizable.

  \item \textbf{Study the decay rate.} The Cathedral's numerical
    experiments show $d_N^2 \sim c/\ln N$ with $c \approx
    1/21.649$. The forward direction is \emph{proved}: the Perron
    chain gives $|M(x)| \leq Cx^{3/4}$ from RH, which yields
    $d_N^2 \leq K/\log N$. Can you prove the conjectured constant?
\end{enumerate}

\subsection*{If You Know Some Physics}

\begin{enumerate}
  \item \textbf{Identify the operator.} The Gram matrix $G_N$ is
    the finite-dimensional shadow of some infinite-dimensional
    operator whose spectral properties encode RH. What is this
    operator? The eigenvalue data ($\lambda_{\min} \sim C/\ln N$,
    $\lambda_{\max} \sim N$, GUE-like statistics) constrain it.

  \item \textbf{Interpret the constant.} The quantum stiffness
    $C \approx 21.649 = 1/(2 + \gamma - \ln 4\pi)$ appears to
    be a physical constant of the ``prime number vacuum.'' Does
    it arise naturally from any known quantum system?
\end{enumerate}

% ================================================================
\section*{What Modern Math Research Actually Looks Like}
% ================================================================

It doesn't look like a genius working alone in a garret.

It looks like this:

\begin{enumerate}
  \item \textbf{Day 1--7}: Explore. Try things. Fail. Archive
    the failures. We produced 128 files of failed attempts. Every
    failure taught us something.

  \item \textbf{Day 8--14}: Find the right framework. For us,
    this was the Parseval Bridge---the realization that frequency-
    domain methods bypass the hardest discrete computations.

  \item \textbf{Day 15--25}: Build. Once the framework is right,
    construction is fast. AI handles the mechanical code generation.
    The human provides architecture. The compiler verifies
    everything.

  \item \textbf{Day 26--30}: Audit. Check every file. Fix every
    docstring. Remove every ghost axiom. Reduce, simplify, clean.
    The Great Audit reduced the codebase to 161 active + 128 archived.
\end{enumerate}

The ratio of exploration to construction is roughly 2:1. Two weeks
of wandering, one week of building. This is normal. If you're not
failing most of the time, you're not exploring hard enough.

% ================================================================
\section*{Why You Should Care}
% ================================================================

The Riemann Hypothesis is not just a mathematical curiosity. It
controls the distribution of prime numbers, which are the atoms
of arithmetic. Every encryption algorithm, every hash function,
every digital signature depends on properties of primes.

But more than that: the RH is a test case for whether humans can
understand the deepest structures in mathematics. If we can prove
it, we learn something profound about the nature of number. If we
can't, we learn something profound about the limits of mathematical
reasoning.

Either way, the attempt is worth making.

% ================================================================
\section*{A Note on Credentials}
% ================================================================

You do not need a PhD to contribute to this project. You do not
need a university appointment. You do not need a research grant.

You need:
\begin{itemize}
  \item A laptop.
  \item Lean~4 (free).
  \item Mathlib (free).
  \item Curiosity.
  \item The willingness to fail 128 times before you succeed once.
\end{itemize}

The compiler does not check your credentials. It checks your proof.

% ================================================================
\section*{The Torch}
% ================================================================

We built the Cathedral. We mapped the territory. We named the
axioms. We preserved the failures. We verified every step.

We could not prove the last two things.

Now it's your turn.

\begin{quote}
\emph{The primes have been singing for 167 years.\\
The Cathedral stands, waiting for its final stones.\\
We built it as far as we could.\\
Two axioms remain.\\
The rest belongs to you.}
\end{quote}

\begin{thebibliography}{99}

\bibitem{baezduarte2003}
L.~B\'aez-Duarte, \emph{A strengthening of the Nyman--Beurling
criterion for the Riemann hypothesis}, Atti Accad. Naz. Lincei
Rend. Lincei Mat. Appl. \textbf{14} (2003), 5--11.

\bibitem{lean4}
L.~de~Moura \emph{et al.}, \emph{The Lean~4 Theorem Prover and
Programming Language}, CADE-28, 2021.

\bibitem{avigad}
J.~Avigad \emph{et al.}, \emph{Mathematics in Lean}, online
textbook, 2023.

\bibitem{buzzard}
K.~Buzzard, \emph{The Natural Number Game},
\url{https://adam.math.hhu.de/\#/g/leanprover-community/nng4}.

\end{thebibliography}

\end{document}
